\documentclass[11pt]{amsart}
\usepackage{amsmath,amssymb,amsthm}
\usepackage[margin=1.1in]{geometry}
\usepackage{newtxtext,newtxmath}
\usepackage{microtype}
\usepackage{booktabs}
\usepackage{tabularx}
\usepackage[colorlinks=true, linkcolor=blue, citecolor=blue, urlcolor=blue]{hyperref}

\newtheorem{proposition}{Proposition}[section]
\newtheorem{remark}[proposition]{Remark}

\setlength{\emergencystretch}{3em}
\raggedbottom

\newcommand{\Q}{\mathbb{Q}}
\newcommand{\F}{\mathbb{F}}
\newcommand{\lean}[1]{\texttt{#1}}
\newcommand{\tagL}{\textsc{[Lean]}}
\newcommand{\tagP}{\textsc{[paper]}}
\newcommand{\tagC}{\textsc{[computation]}}

\begin{document}

\title[PR589 hosting and realizability]{Hosting and realizability in the selected PR589 bit supplier}

\author{Dan Brown$^*$}
\email{jdb1729@gmail.com}
\thanks{$^*$Author of \emph{The Da Vinci Code}.}
\author{Kenny Daniel}
\email{platypii@gmail.com}

\date{October 18, 2026}

\begin{abstract}
The selected bit supplier of the community bound on integer multiplication~\cite{CS}, PR589 on
PR555, prices one interchange node at the saving $a_b = 10/5769$. Its accounting makes every
rank-$r$ frame change of its word one child interchange of width $r b$, and its routes and charts
free. We record what happens when this node is ported to the recursive interchange machine that is
formalized in Lean~\cite{IMB}. (A) In one invocation a family hosts up to five helpers. Upstream
separates them by classes of an orthogonal-group cover. A machine with a fixed role count can
instead split each family into five row classes, which leaves the occurrence moment unchanged. The
role-level node theorem of this scheme is checked in Lean.
(B) The machine's cheap operations are lower triangular, so a frame change whose projector
difference is not diagonal is not one child, and the routes and charts are coordinate changes that
cost interchanges. Oblique flag frames, with each frame change split into runs of consecutive pivots, turn the
frame changes into cheap conjugates of children.
The routes remain: in a stated machine model, which charges sub-field transpositions as children of
the same recursion, they raise the occurrence moment of upstream's compressed geometry above one,
and no saving remains. The saving needs the compression, since an uncompressed block geometry
loses its only free coordinate. These are arguments on paper with computation, not Lean theorems.
The no-go holds for architectures without class-indexed covers. Upstream's later packages pay for
the routes with a cover indexed by $\mathrm{GL}_{64}(\mathbb{Z}/q^w)$, a different machine whose
completion is open by upstream's own account. The finite certificates of the supplier, re-checked in Lean, agree with upstream.
\end{abstract}

\maketitle

\section{Introduction}

The verified multiplier of~\cite{IMB} runs in time $O\bigl(n(\log n)^{1-\kappa}\bigr)$ on Mathlib's
finite multi-stack Turing machines, with $\kappa = 651 \cdot 2^{-39}$. Upstream's selected record~\cite{CS} combines a complex
supplier (T16) with a bit supplier (PR589, tuned by DreamingOfClouds on eumemic's PR555) and states
the conditional exponent $\kappa = 0.001672660197076918$. Its README lists the open parts: ``All-size/compiler,
common-label/mutable-source realization, whole-call hosting, precision/tape, exact recovery and
analytic obligations remain conditional.''

The Lean development is porting both selected suppliers into its machine. Under upstream's
accounting, the two together would move the cap on $\kappa$ from about $2^{-29.6}$ to about
$2^{-10.3}$; either alone leaves it near $2^{-29.4}$. This note records what the port found on the bit side.

Each claim carries a tag. \tagL{} marks a theorem checked in~\cite{IMB}, under the name given, with
only Lean's standard axioms. \tagP{} marks a written argument, here or in the design notes of the
development, that is not formalized. \tagC{} marks a number computed by a script from upstream's
pinned data. Upstream sources are pinned at the scientific base \texttt{0fdd7d82} and, for later
packages, at \texttt{463113d6}.

\section{The supplier and its bill}

The parent of a bit interchange node holds records whose addresses are pairs $(x,y)$, with $x$ the
$H$ field and $y$ the $D$ field, each of $64$ coordinates; the interchange exchanges them. For an
idempotent $U$ write
\[
  D_U(x, y) = \bigl((1-U)x + Uy,\; Ux + (1-U)y\bigr),
\]
so the parent interchange is $D_1$ (\lean{exchange} in \lean{Nine/PortFrame}). The supplier's word
has $5583$ instructions on $661$ registers. Each register carries a frame: a label $V \subseteq
\Q^8$ spanned by $E_8$ lines, with orthogonal projector $P_V$. A MOVE takes a register from $P_V$ to
$P_{V'}$ with $V \subseteq V'$, and \lean{exchange\_change} \tagL{} gives
$D_{P_V} D_{P_{V'}} = D_{P_{V'} - P_V}$. Its rank is $\operatorname{rank}(P_{V'} - P_V)$.

The node runs the word in nine stages over four pairs of data banks, in $32$ replicas, on $111$
live ports, with $439$ helpers per invocation hosted in idle families. It has $W = 28416$ live roles.
The accounting makes every rank-$r$ MOVE of every invocation, and every hosted idle remainder of
rank $r$, one child interchange of width $r b$, where $b = w/64$ for a parent of width $w$. The
children of one node then have the rank histogram
\begin{center}
\begin{tabular}{lcccccccc}
\toprule
rank & 1 & 2 & 3 & 4 & 5 & 6 & 7 & 42\\
count & 381024 & 400416 & 22752 & 26880 & 29376 & 17376 & 26784 & 288\\
\bottomrule
\end{tabular}
\end{center}
which is \lean{bitHistogram}, equal to upstream's profile. \tagL{}
\begin{itemize}
\item \lean{nodeChildren\_realizes}: the node's children, in schedule order, realize
  \lean{bitHistogram}.
\item \lean{bit\_bill\_le}: the occurrence moment $\frac{1}{W}\sum_{\text{children}} (r/64)^{1-a_b}$
  plus its surcharge is at most $1$, at $a_b = 10/5769 \approx 1.7334 \cdot 10^{-3}$.
\item \lean{bit\_variable\_closed}: a node whose children have sizes at most $r/64$ of the parent
  (the premise \lean{hsize}) and weight $1/W$ costs $O(e^\lambda)$ for every
  $\lambda > 1 - a_b$.
\end{itemize}
These results take the accounting as given. Its two assumptions are the subject of this note:
a rank-$r$ frame change is one child of width $r b$, and routes and charts are free.

\section{Hole A: hosting}

\subsection{Five helpers in one family}
In one invocation a family of an inactive pair can host up to five helpers. They lie on slots of a
single idle (\lean{same\_pair\_windows} \tagL{}) at offsets $8j$ with $j < 5$ (\lean{slot\_fits}
\tagL{}). Upstream's hosting audit~\cite{HA} counts the same maximum of five. In upstream's first
example~\cite{FH}, at stage~2, replica~0, pair~0, helpers $232$ through $236$ share family $666$, on
the supports $1..8$, $9..16$, $17..24$, $25..32$ and $33..40$ of idle child~4's $42$-dimensional
pending projector. The helpers are separate registers with possibly correlated contents, so they
need separate storage.

\subsection{Upstream's cover}
Upstream separates them with a cover~\cite[\S\S2,4]{FH}. A physical stream is indexed by a pair
$(F, g)$, with $F$ a family and $g$ an element of the full orthogonal group $G$ of the form
$I_{64}/2$ over a finite field, and a cover sweep visits every element of $G$ once. At invocation
$d$ the helper on slot $S$ is owned by $(F, d\sigma_S^{-1})$, so helpers that share $F$ have
different classes. Upstream's later row allocations carry this group in the row stock, through the
factor $\lvert \mathrm{GL}_{64}(\F_q)\rvert\, q^{4096(w-1)}$~\cite{LRS}. The machine
of~\cite{IMB} has a fixed role count, and its row divisor is a power of that count
(\lean{RecursiveRowDivisor}). It has no cover classes.

\subsection{Five row classes}
Split each family's rows into five classes. Run each invocation once per class $d \in
\mathbb{Z}/5$, and give the helper on the slot at offset $8j$ the class $(d + j) \bmod 5$ of its
host family (\lean{hostClass}).
\begin{itemize}
\item \tagL{} \lean{hostClass\_distinct}: two distinct slots of one stage, pair and copy receive
  distinct classes for every $d$. \lean{hostClass\_bijective}: as $d$ runs through $\mathbb{Z}/5$,
  each slot visits each class of its family once. \lean{card\_role}: the node has $5W = 142080$
  roles.
\item \tagL{} One class run, on these roles: both sides of the active pair change between their
  stage frames, and a hosted helper's role changes by the exchange on its slot
  (\lean{invoke\_side\_zero}, \lean{invoke\_side\_one}, \lean{invoke\_host}). All $160$ class runs
  of a stage: two runs host orthogonal slots on a role (\lean{hostSpan\_orthogonal}), each residual
  of the stage arrives on its role exactly once (\lean{hostSpan\_sum}), and each role of another
  pair changes by the exchange on the sum of the slots its family hosts at that stage
  (\lean{stagePiece\_inactive}).
\item \tagL{} The node: \lean{nodeProgram} has $29$ pieces, and on arbitrary dirty streams it
  changes every role by $D_1$ (\lean{runNode\_interchange}). Every child appears five times, with
  weight $1/(5W)$: the child list realizes five times \lean{bitHistogram}, and the occurrence
  moment $\frac{1}{5W}\cdot 5\sum (r/64)^\lambda$ is at most $1$ at $\lambda = 1 - a_b$
  (\lean{five\_class\_node}). \lean{five\_variable\_closed} is the closed recurrence at
  weight $1/(5W)$, with the per-node cost as a hypothesis.
\end{itemize}
These semantics are at role level. Streams are functions of rational address pairs; routes,
charts and frame changes act on them as exact operators; and each frame change is priced as one
child of its rank. The scheme is sound at that level. Its realization on tape depends on hole~B.

\section{Hole B: realizability}

\subsection{The machine model}\label{sec:model}
The arguments of this section use the following model of a node of the recursive interchange
machine, on volume $V$, with each address field split into sub-fields of $b$ digits stored in
$H$-before-$D$ order.
\begin{enumerate}
\item[(1)] The nonrecursive work of a node costs $O(V)$, with a constant independent of the
  width.
\item[(2)] Linear passes are: scalings; shifts whose control sub-field is more significant than
  the target in the $H$-before-$D$ layout; cross shears $D \mathrel{+}= M H$; XORs and copies of
  role arrays; and role relabeling. On the address coordinates of one role these act by lower
  triangular matrices. In~\cite{IMB} the scalings and shifts are \lean{FlatCoordinateSchedule.Op},
  whose shift control precedes its target. The shear programs of the manuscript~\cite{OAI} have the same form:
  Lemma~4.2, \lean{exists\_prime\_shear} \tagL{}, uses exactly $\operatorname{rank} A$
  interchanges and otherwise lower triangular transformations, and every transformation of the
  generated Shared50 program is lower triangular (\lean{programs\_triangular} \tagL{}).
\item[(3)] Moving data from a less significant into a more significant sub-field, as in upper
  transvections and sub-field transpositions, is charged as an interchange child, since a
  transposition of two sub-fields of width $b$ is itself an interchange instance. A child of
  rank $r$ swaps $r$ consecutive $H$ sub-fields with $r$ consecutive $D$ sub-fields, at width
  $r b$ (\lean{childR} \tagL{}).
\item[(4)] Upstream's cover model sidesteps (3), and upstream lists the all-size compiler as an
  inherited assumption.
\end{enumerate}

\subsection{Frame changes}
\begin{proposition}[\tagP]\label{prop:triangular}
Let $A$ and $B$ be lower triangular in the $H$-before-$D$ order, and let $E$ be a diagonal
$0/1$ matrix. The upper-right block of $A D_E B$ is lower triangular. Hence, if $U$ is symmetric
and not diagonal, $D_U$ is not of the form $A D_E B$.
\end{proposition}

\begin{proof}
Write $A = \begin{pmatrix} A_1 & 0\\ A_3 & A_4\end{pmatrix}$ and $B = \begin{pmatrix} B_1 &
0\\ B_3 & B_4\end{pmatrix}$, with $A_1, A_4, B_1, B_4$ lower triangular. Since $D_E =
\begin{pmatrix} 1-E & E\\ E & 1-E\end{pmatrix}$, the upper-right block of $A D_E B$ is $A_1 E B_4$,
a product of lower triangular matrices. The upper-right block of $D_U$ is $U$, and a symmetric lower
triangular matrix is diagonal.
\end{proof}

So with orthogonal frames, a MOVE whose projector difference $P_{V'} - P_V$ is not diagonal is not
one child between linear passes. In the word's local coordinates, $2787$ of its $2911$ paid MOVEs
(phase incidences of positive rank) have a non-diagonal difference \tagC.

\tagP{} Oblique flag frames remove this obstacle. Give each label $V$ the idempotent $\Pi_V$ with
range $V$ whose kernel is spanned by the coordinate vectors outside the pivot set $S_V$ of $V$, the
set of leading positions of nonzero vectors of $V$. From a reduced echelon basis, $\Pi_V$ is lower
triangular with diagonal $1_{S_V}$. For $V \subseteq V'$ one has $S_V \subseteq S_{V'}$ and
$\Pi_V \Pi_{V'} = \Pi_{V'} \Pi_V = \Pi_V$, so \lean{exchange\_change} applies with these frames,
and the frame change is $D_U$ with $U = \Pi_{V'} - \Pi_V$, of rank $\dim V' - \dim V$. The
idempotent $U$ is lower triangular with diagonal $E = 1_{S_{V'} \setminus S_V}$, so $U = C E
C^{-1}$ for a lower triangular $C$, and $D_U = \operatorname{diag}(C,C)\, D_E\,
\operatorname{diag}(C,C)^{-1}$ (\lean{exchange\_conj} \tagL{}). The frame change is a cheap
conjugate of the swap of its new pivot coordinates. The Lean frame identities of \S6 use orthogonal
frames; their oblique counterparts are not formalized.

\tagC{} The new pivots need not be consecutive, and \lean{childR} needs consecutive sub-fields.
Over the $40320$ orders of the local axes, the order $(3,0,2,4,6,5,7,1)$ gives consecutive new
pivots for $2281$ of the $2911$ paid MOVEs, which carry $3458$ of the rank mass $5066$ ($68\%$).
Splitting each MOVE into one child per maximal run of consecutive new pivots gives the run
histogram $\{1{:}2443,\ 2{:}979,\ 3{:}86,\ 4{:}35,\ 5{:}35,\ 6{:}6,\ 7{:}8\}$, $3592$ runs. Each run
is a cheap conjugate of a \lean{childR} child, since the exchanges on orthogonal diagonal
projectors multiply (\lean{exchange\_mul\_of\_orthogonal} \tagL{}).

\subsection{Routes in the compressed geometry}
The stage route of a port with line $q$ is $\rho = \begin{pmatrix} P & A\\ A^{\mathsf T} & 0
\end{pmatrix}$ on the eight local axes and one exterior block of seven, with $P = qq^{\mathsf T}/4$.
It is a symmetric involution (\lean{route\_transpose}, \lean{route\_mul\_self} \tagL{}), with
entries in $\{0, \pm\frac14, \pm\frac12, \pm 1\}$, and a lifted MOVE is the route conjugate of the
local one, extended by zero (\lean{lift\_sub} \tagL{}). The slot charts embed a helper's eight
local axes isometrically onto its slot (\lean{slotChart\_orthonormal} \tagL{}). Routes and charts
are orthogonal and not triangular. In upstream's cover model they are free: $g\sigma = d$ for every
role, so ``a chart-routed ADD is exactly a common native address ADD''~\cite[\S2]{FH}. On a tape
they are coordinate changes, that is, permutations of the address space applied to stored records.

In the model of \S\ref{sec:model}, with route transpositions executed as children of this
recursion, the routes cost more than the saving.
\begin{itemize}
\item \tagC{} Every invertible matrix factors as $L_1 \pi L_2$ with $L_1, L_2$ lower triangular and
  $\pi$ a unique permutation (Bruhat). With the local axes in the order $(3,0,2,4,6,5,7,1)$ and the
  exterior block after them, the permutation of each of the $111$ port routes has $7$
  transpositions on each field.
\item \tagP{} The routes must depend on the port: a route common to all ports would map all $111$
  lines to one axis.
\item \tagP{} By (3), executing a route through its Bruhat form costs sub-field transpositions,
  each an interchange child. The feasibility study counts about $4.5$ rank-seven interchanges, of
  width $7b$, per role per node.
\item \tagP{} At $\lambda = 1$ the occurrence moment is the rank mass over $64W$. Without routes
  the mass per role is $1808352/W \approx 63.6385$, so the moment is $63.6385/64 \approx 0.99435$,
  a slack of about $0.0056$. The routes add about $4.5\,(7/64)^{\lambda} \approx 0.49$ to the moment
  near $\lambda = 1$, so it exceeds $1$ for every $\lambda \le 1$. No saving remains.
\item \tagC{} Without routes, splitting each MOVE into consecutive runs would give the saving
  $1.6407 \cdot 10^{-3}$, and rank-one children $1.362 \cdot 10^{-3}$. Both still need the routes of
  the compressed geometry, so neither escapes the count above.
\end{itemize}

\subsection{Why the saving needs the compression}
\tagP{} A bit saving needs rank mass below $mW$ at $\lambda = 1$, for parent width $m$. The only
free coordinate of PR589 is the shared port line. It is worth $1$ per data column, against a COPY
overhead of $L \cdot \mathrm{CST}/(2KV) = 9 \cdot 63/888 \approx 0.638$, where $\mathrm{CST} = 63$
is the sum of the nine COPY ranks. That leaves about $0.362$ per role, and
$W \cdot 0.362 \approx 10272$ is upstream's deficit $T(2KV - L \cdot \mathrm{CST})$ at $T = 32$,
$K = 4$, $V = 111$, $L = 9$~\cite{TW}.

In an uncompressed block geometry, with each stage running the word on its own block of eight axes
and parent width $m = 72$, no routes or charts are needed. Without compression, the
asymmetry between the two sides leaves the stage window after the first stage. Moving the line to a
port-independent axis $u$ makes the bank matrix a scalar multiple of the identity, so the free
coordinate cancels \tagC. The mass per role is then $72.638 > 72$, and the moment at $\lambda = 1$
is $1.0089 > 1$. The uncompressed geometry has no saving.

\subsection{Upstream's later treatment of the routes}
After the base, upstream merged packages that price the routes explicitly, at
\texttt{463113d6}. \texttt{selected-bit-stream-routes}~\cite{SR} compiles the $999$ data and
$124$ hosted route types into exact rational column words and charges up to $3520$ ordinary
interchange calls per route action. \texttt{selected-bit-recursive-caller}~\cite{RCa} states that
``its routing and leaf swaps still cost $O(wV)$''.
\texttt{selected-bit-fast-primitives}~\cite{FP} replaces both ordinary-interchange sites with a
previously constructed native interchange and claims a local toll $O_q\bigl(V(n +
w^{1-a_1})\bigr)$ with $a_1 = 169711/10^8$. \texttt{selected-bit-recursive-closure}~\cite{RCl}
claims ordinary binary chunk interchange in $O(V E^{1-a})$ with $a = 0.00169711$, in its fixed-tape
complete-stream model. Their authors present them as written arguments without independent or
proof-assistant verification.

\tagP{} Their toll rests on a class-indexed cover. Every stream is indexed by $g \in
\mathrm{GL}_{64}(\mathbb{Z}/q^w)$, written $g = g_0 + qG$ with a low class $g_0 \in
\mathrm{GL}_{64}(\F_q)$ and $4096$ high fields $G_{ij}$ of $w-1$ digits, and each packet is ordered
as its low-class label, then the high fields, then the payload~\cite[PROOF.md \S1]{SR}. A route is
the right action $g \mapsto g\sigma$ on this label of fixed size: a block permutation of the low
classes, plus carry-aware updates of the high fields through the previous-stage interchange at width
$w-1$ (``B1, ordered ADD, B1''~\cite[PROOF.md \S4]{FP}). A route therefore costs $O_q(V w^p)$ for a
fixed exponent $p$, independent of the payload width, because it never moves payload address
digits.

\tagP{} In the model of \S\ref{sec:model} roles carry no class label, so a route acts on the
payload's own $64$ sub-fields of $n/64$ atoms each. Whether it runs as $q$-adic sub-field arithmetic
or as atom-wise elementary operations, its cost per node is about $V n w^p$: a product of the two
widths where upstream's toll has a sum, superlinear in the payload width.

\tagP{} The cover costs volume. The stock carries a factor $\lvert
\mathrm{GL}_{64}(\mathbb{Z}/q^w)\rvert$, constant for fixed $q$ and $w$, with $\lvert
\mathrm{GL}_{64}(\F_q)\rvert > 2^{327680}$ at $q > 2^{80}$, and children need their own classes,
through the stock and a fresh-class contract. The saving also needs $w$ to grow: at fixed $w$ the
time per unit volume is $\Theta(n)$. It comes from balancing $n \approx e^{\theta}$ against $w
\approx e^{1-\theta}$, which gives $a_{k+1} = c(1-c) + c\,a_k$ with $c = 17/10000$~\cite{FP}.

\tagC{} If that architecture were completed, one bootstrap from the Shared50 interchange would give
$a_b \approx c(1-c) \approx 1.697 \cdot 10^{-3}$, and with T16, $\kappa \approx 2^{-10.82}$ with
rank-one T16 children or $\kappa \approx 2^{-10.44}$ with variable-rank ones.

So the no-go of \S4.3 holds for architectures without class-indexed covers. Upstream's cover
architecture is a different machine, and its completion is open by its own account: ``We retain
the data-dependent recursive compiler, weighted calls, correlated-helper domain bridge and COPY
lifetime obligations. [\dots] The selected result remains conditional.''~\cite{SR}

\section{Paths forward}

Besides completing upstream's cover architecture (\S4.5), three paths remain.
\begin{enumerate}
\item[(i)] A bit word whose sources share a port-independent line, so that compression needs no
  routes.
\item[(ii)] A primitive that makes the needed sub-field moves linear, against which assumption (3)
  of \S\ref{sec:model} argues.
\item[(iii)] A different bit supplier.
\end{enumerate}
The hosting scheme of hole~A and the oblique frames of \S4.2 apply to any compressed geometry
with these hosting windows, so they carry over to path~(i).

\section{What holds}

The finite certificates of the supplier are re-checked in Lean and agree with upstream. \tagL{}
\begin{itemize}
\item Local MOVE profile: every phase incidence is a nested frame pair whose rank is the certified
  dimension difference (\lean{phase\_ranks}, \lean{phase\_histogram}); final alignment pays no rank
  (\lean{final\_ranks}); the nine COPY ranks are all $7$ (\lean{copyRanks\_eq}); the local paid
  histogram is upstream's $\{1{:}1323,\ 2{:}1342,\ 3{:}79,\ 4{:}44,\ 5{:}102,\ 6{:}11,\ 7{:}19\}$
  (\lean{localPaid\_histogram}).
\item Stage geometry: the $111$ port certificates repeat upstream's reflection walk
  (\lean{ports\_checked}); each column leaves a stage in the frame in which it enters the next
  (\lean{stage\_source}); after nine stages and the exchange, every column's frame is $D_1$ after its
  initial frame (\lean{node\_exchange\_source}); lifts are route conjugates (\lean{lift\_sub}).
\item Windows: upstream's nine-stage windows (\lean{windows\_four}), and every hosted stage lies in
  its idle's window (\lean{selection\_in\_windows}).
\item Occurrence link and bill: \lean{nodeChildren\_realizes}, \lean{bit\_bill\_le}, and the
  recurrence \lean{bit\_variable\_closed}.
\end{itemize}
The occurrence link and the recurrence stand, given the accounting. Hole~B asks whether a machine
realizes the accounting.

The complex side has no analogue of either hole with rank-one children. Kenny Daniel's T16 node
program runs on a fixed set of $14336$ data families and eight private centres, with no cover
classes. Every family ends with $C_{\mathrm{all}}$ applied to its input (\lean{nodeProgram\_cAll}
\tagL{}), its child calls are $1041408$ rank-one odd kernels (\lean{nodeProgram\_children}
\tagL{}), and in every column such a kernel is the layer machine's factor (\lean{colRep\_kernel}
\tagL{}).

\section{Status}

\begin{center}
\small
\begin{tabularx}{\textwidth}{@{}Xl@{}}
\toprule
item & status\\
\midrule
Profile, stage geometry, windows, occurrence link, bill, recurrence & Lean\\
Five classes: class facts, one class run, all class runs of a stage & Lean\\
Five classes: node program, node theorem, moment at weight $1/(5W)$ & Lean\\
A MOVE with non-diagonal difference is not one child (Prop.~\ref{prop:triangular}) & paper\\
$2787$ of $2911$ paid MOVEs have a non-diagonal difference & computation\\
Oblique flag frames make frame changes cheap conjugates of children & paper\\
Pivot runs at the best axis order & computation\\
Bruhat counts of the $111$ routes & computation\\
Compressed geometry: no saving in the model of \S\ref{sec:model} & paper and computation\\
Uncompressed block geometry at $m = 72$: no saving & paper and computation\\
Cheap routes need a class-indexed cover (\S4.5) & paper\\
Upstream's cover architecture & open by upstream's account\\
\bottomrule
\end{tabularx}
\end{center}

\begin{thebibliography}{9}\raggedright

\bibitem{CS} D. Colkitt (maintainer) and contributors, \texttt{integer-mult-bounds}: community
research on integer multiplication below $n \log n$, base \texttt{0fdd7d82}, head
\texttt{463113d6}, \url{https://github.com/CrocSwap/integer-mult-bounds}.

\bibitem{FH} \texttt{research/pr589-five-helper-hosting}, PROOF.md, in~\cite{CS} at
\texttt{0fdd7d82},
\url{https://github.com/CrocSwap/integer-mult-bounds/tree/0fdd7d82de03e43c9e5a1000dfbcd380d1b687d2/research/pr589-five-helper-hosting}.

\bibitem{FP} \texttt{research/selected-bit-fast-primitives}, in~\cite{CS} at \texttt{463113d6},
\url{https://github.com/CrocSwap/integer-mult-bounds/tree/463113d675af7410ac91d53a4fed51ec6494eb9a/research/selected-bit-fast-primitives}.

\bibitem{HA} \texttt{research/selected-bit-hosting-audit}, in~\cite{CS} at \texttt{0fdd7d82},
\url{https://github.com/CrocSwap/integer-mult-bounds/tree/0fdd7d82de03e43c9e5a1000dfbcd380d1b687d2/research/selected-bit-hosting-audit}.

\bibitem{IMB} K. Daniel, \texttt{integer-mult-bounds-lean}: integer multiplication
below $n \log n$ in Lean on Mathlib's finite Turing machines,
\url{https://github.com/platypii/integer-mult-bounds-lean}.

\bibitem{LRS} \texttt{research/selected-bit-local-ring-sweeps}, in~\cite{CS} at
\texttt{463113d6},
\url{https://github.com/CrocSwap/integer-mult-bounds/tree/463113d675af7410ac91d53a4fed51ec6494eb9a/research/selected-bit-local-ring-sweeps}.

\bibitem{OAI} OpenAI,
\emph{Integer multiplication below $n \log n$}, OpenAI Math Release preprint, September 23, 2026,
\url{https://github.com/openai/math/blob/main/preprints/Integer-multiplication-below-n-log-n-September-23-2026/paper.pdf}.

\bibitem{RCa} \texttt{research/selected-bit-recursive-caller}, in~\cite{CS} at
\texttt{463113d6},
\url{https://github.com/CrocSwap/integer-mult-bounds/tree/463113d675af7410ac91d53a4fed51ec6494eb9a/research/selected-bit-recursive-caller}.

\bibitem{RCl} \texttt{research/selected-bit-recursive-closure}, in~\cite{CS} at
\texttt{463113d6},
\url{https://github.com/CrocSwap/integer-mult-bounds/tree/463113d675af7410ac91d53a4fed51ec6494eb9a/research/selected-bit-recursive-closure}.

\bibitem{SR} \texttt{research/selected-bit-stream-routes}, in~\cite{CS} at \texttt{463113d6},
\url{https://github.com/CrocSwap/integer-mult-bounds/tree/463113d675af7410ac91d53a4fed51ec6494eb9a/research/selected-bit-stream-routes}.

\bibitem{TW} \texttt{research/e8-bit-p111-tuned/code/adapt9p.py}, in~\cite{CS} at
\texttt{0fdd7d82},
\url{https://github.com/CrocSwap/integer-mult-bounds/blob/0fdd7d82de03e43c9e5a1000dfbcd380d1b687d2/research/e8-bit-p111-tuned/code/adapt9p.py}.

\end{thebibliography}

\end{document}
